\documentclass[12pt,a4paper]{article}
\usepackage[utf8]{inputenc}
\usepackage[T1]{fontenc}
\usepackage{geometry}
\usepackage{hyperref}
\usepackage{listings}
\usepackage{xcolor}
\usepackage{graphicx}
\usepackage{float}

% Margins
\geometry{margin=1in}

% Code highlighting style
\definecolor{codegray}{rgb}{0.95,0.95,0.95}
\definecolor{commentgreen}{rgb}{0,0.6,0}
\definecolor{keywordblue}{rgb}{0,0,0.6}

\lstdefinestyle{mystyle}{
    backgroundcolor=\color{codegray},
    commentstyle=\color{commentgreen},
    keywordstyle=\color{keywordblue},
    basicstyle=\ttfamily\small,
    breakatwhitespace=false,
    breaklines=true,
    captionpos=b,
    keepspaces=true,
    showspaces=false,
    showstringspaces=false,
    showtabs=false,
    tabsize=2,
    frame=single
}
\lstset{style=mystyle}

\title{\textbf{CoFound Application}\ \large Individual Project, 2025 - 2026 Instruction Manual}
\author{Rus-Gheorghiu Lucas-Michael \\ Faculty of Computer Science \\ Year 2, Group 5}
\date{\today}

\begin{document}

\maketitle

\tableofcontents
\newpage

\section{Introduction}
\textbf{CoFound} is a dedicated collaboration platform designed to bridge the gap between visionary startup founders and skilled professionals looking to join a team. The application is intended for students, developers, and entrepreneurs who wish to find the perfect co-founder or project partner.

The main purpose is to facilitate meaningful connections through detailed profiles, skill matching, and project-based communication.

The primary functionalities of the application include:
\begin{itemize}
    \item \textbf{User Profiling & Skills} – Users can create detailed profiles with biographies, professional links (LinkedIn, GitHub), and a list of technical skills.
    \item \textbf{Project Management} – Founders can create project listings detailing their idea, required roles, and team size.
    \item \textbf{Application System} – Interested candidates can apply to join projects, and owners can review, accept, or reject these applications.
    \item \textbf{Real-time Communication} – Integrated chat functionality allows for direct communication between project members.
\end{itemize}

By using this application, users will benefit from a streamlined recruiting process and a centralized hub for managing startup collaborations.

\section{Design}
The application architecture is robust and modern, utilizing the following technologies:

\begin{itemize}
    \item \textbf{Backend:} Java Spring Boot (v3.2.5) providing a RESTful API.
    \item \textbf{Frontend:} React.js (v18) with Bootstrap 5 for a responsive user interface.
    \item \textbf{Database:} H2 Database (for local development persistence) / PostgreSQL (production ready).
    \item \textbf{Security:} JWT (JSON Web Tokens) for stateless authentication and role-based access control.
    \item \textbf{Build Tools:} Maven (Backend) and Vite (Frontend).
\end{itemize}

\subsection{Database Schema}
The data model revolves around three core entities:
\begin{itemize}
    \item \textbf{User:} Stores authentication credentials, profile data, and contact information.
    \item \textbf{Project:} Contains details about the startup idea, required team size, and status. It has a One-to-Many relationship with Applications.
    \item \textbf{Skill:} A Many-to-Many relationship with Users, allowing for tagging and searching candidates based on expertise (e.g., Java, Python, Marketing).
\end{itemize}

\subsection{Security Implementation}
Security is a priority for CoFound.
\begin{itemize}
    \item \textbf{Password Hashing:} All user passwords are encrypted using BCrypt before storage.
    \item \textbf{Stateless Auth:} The API uses JWTs. Upon login, the server issues a token which the client presents in the Authorization header for subsequent requests.
    \item \textbf{Role-Based Access:} Endpoints are protected so that only project owners can modify their projects or view private applications.
\end{itemize}

\section{Diagrams}
\textit{(You can insert your UML diagrams here, such as Class or Use Case diagrams, to illustrate the system structure.)}

% Example placeholder for an image
% \begin{figure}[H]
%     \centering
%     \includegraphics[width=0.8\textwidth]{diagram.png}
%     \caption{System Architecture Diagram}
% \end{figure}

\section{Implementation}
To download and continue the development of the application, the client can follow these steps:

\subsection{1. Cloning the Repository}
\begin{lstlisting}[language=bash]
git clone https://github.com/lucas-rus/cofound-app.git
cd cofound-app
\end{lstlisting}

\subsection{2. Configuration (Critical)}
For security, you must configure your own email credentials.
\begin{enumerate}
    \item Navigate to \texttt{src/main/resources/application.properties}.
    \item Open the file and update the mail settings with your own Gmail App Password:
\end{enumerate}
\begin{lstlisting}[language=bash]
spring.mail.username=YOUR_EMAIL@gmail.com
spring.mail.password=YOUR_APP_PASSWORD
\end{lstlisting}

\subsection{3. Starting the Backend}
Open a terminal in the project root:
\begin{lstlisting}[language=bash]
# Ensure Java 17 is used
./start-backend.sh
# OR manually:
mvn spring-boot:run
\end{lstlisting}
The server will start on \texttt{http://localhost:8080}.

\subsection{4. Starting the Frontend}
Open a second terminal in the \texttt{frontend} folder:
\begin{lstlisting}[language=bash]
cd frontend
npm install  # First time only
npm run dev
\end{lstlisting}
The client will launch on \texttt{http://localhost:5173}.

\section{Usage}

\subsection{1. Authentication and Registration}
Users can securely create an account by providing an email and password. A verification link is sent via email to ensure authenticity. Once verified, users log in using JWT-based authentication.

\subsection{2. Interface Navigation}
The GUI is designed for intuitive navigation:
\begin{itemize}
    \item \textbf{Dashboard:} The landing page displaying all available projects.
    \item \textbf{Navigation Bar:} Provides quick access to "Create Project", "My Applications", and the User Profile.
    \item \textbf{Profile Page:} A centralized hub for managing personal data and settings.
\end{itemize}

\subsection{3. Using Main Functions}

\subsubsection{Profile Management}
Users can personalize their presence on the platform:
\begin{itemize}
    \item \textbf{Data Entry:} Add a bio, upload a CV/Resume, and set a Profile Picture.
    \item \textbf{Skill Tagging:} Users can add skills (e.g., "Java", "Marketing") which help in matching with projects.
    \item \textbf{Account Control:} Users have full control to permanently delete their account and associated data via a secure confirmation modal.
\end{itemize}

\subsubsection{Project Creation & Workflow}
Founders can post projects with specific requirements.
\begin{itemize}
    \item \textbf{Creation:} Define a title, description, and required skills.
    \item \textbf{Application:} Users view project details and click "Apply".
    \item \textbf{Status Tracking:} The status of an application transitions from \texttt{PENDING} $\rightarrow$ \texttt{ACCEPTED} or \texttt{REJECTED}.
    \item \textbf{Review:} Project owners receive applications and can "Accept" or "Reject" them via the project dashboard.
\end{itemize}

\section{Usage Examples}
\textit{(Screenshots of the application can be placed here.)}

\subsection{Example 1: Creating a Project}
The "Create Project" form validates input and ensures all necessary details like "Team Size" and "Category" are filled before submission.

\subsection{Example 2: Account Deletion}
To delete an account, the user must navigate to the Profile page, scroll to the bottom, and click "Delete Account". A modal appears requiring the user to type "delete" to confirm, preventing accidental data loss.

\section{Conclusions and Next Steps}
This application provides a robust solution for the startup ecosystem, allowing efficient team formation. Users benefit from an optimized and secure experience.

The application can be improved by:
\begin{itemize}
    \item \textbf{Mobile Application:} Developing a native mobile app (React Native or Flutter) to allow users to manage applications on the go.
    \item \textbf{Calendar Integration:} Integrating tools like Google Calendar or Calendly to schedule interviews directly within the chat.
    \item \textbf{Premium Features:} Implementing a monetization model where founders can pay to highlight their projects for better visibility.
    \item \textbf{Advanced Analytics:} Providing a dashboard for project owners to track views, clicks, and conversion rates of their project posts.
\end{itemize}

\end{document}
